\section{Robustness}
\label{sec:robustness}

% Cada teste responde a uma ameaca nomeada. WCB e permutacao: poucos clusters
% (tab07, na secao de resultados). Placebo de antecipacao: datacao e
% pre-tendencias. Honest DiD: violacoes limitadas de tendencias paralelas.

The anticipation placebo shifts each treatment date twelve months earlier
and is estimated on the true pre-period only. For cabotage $T_2$ it returns
\PlaceboCabTdoisBeta\ ($p=\PlaceboCabTdoisPwcb$): the apparent reduction in
$T_2$ predates adoption and is therefore a pre-trend, not an effect. The
$T_4$ placebo is not distinguishable from zero ($p=\PlaceboCabTquatroPwcb$).

\input{tables/tab09_placebo}

\input{tables/tab12_lee_bounds}

\begin{figure}[htbp]
\centering
\includegraphics[width=\textwidth]{fig14_estimates.pdf}
\caption{Static difference-in-differences estimates by sample (design A). Dots and thin
lines: point estimate and 95\% cluster-robust interval; wide shaded bars: Lee
bounds from \Cref{tab:lee} where computed; $p$ (right): wild cluster
bootstrap $p$-value (Webb weights). Source: ANTAQ waterway statistics, own elaboration.}
\label{fig:estimates}
\end{figure}

\input{tables/tab13_pretrend}

\begin{figure}[htbp]
\centering
\includegraphics[width=\textwidth]{fig15_honestdid.pdf}
\caption{Robust confidence intervals for the average cabotage effect (event quarters 0--4)
under the smoothness restriction of \citet{rambachan2023}. $M=0$ allows the pre-adoption
trend to continue linearly after adoption; larger $M$ allows curvature. ``Original'' is
the conventional 95\% interval. Design A, Sun--Abraham event study in quarters with
port, month and navigation fixed effects and cargo controls; standard errors clustered
by port. Source: ANTAQ waterway statistics, own elaboration.}
\label{fig:honestdid}
\end{figure}

\PROSA

\PROSA

\FloatBarrier
